\documentclass[10pt]{mypackage}

\usepackage[uprightgreeks,light,noDcommand]{kpfonts}
\renewcommand{\coloneq}{\coloneqq}
\renewcommand{\eqcolon}{\eqqcolon}
%\renewcommand{\mathbb}{\mathds}
%\usepackage{homework}
\usepackage{notes}

\usepackage[ backend=bibtex, style = alphabetic, sorting=ynt ]{biblatex}
\addbibresource{all_references.bib}

\usepackage{parskip}

\fancyhf{}
\fancyhead[R]{Avinash Iyer}
\fancyhead[L]{Sofic and Hyperlinear Groups}
\fancyfoot[C]{\thepage}

\setcounter{secnumdepth}{0}

\begin{document}
\RaggedRight
\section{Ultrapowers}%
If $\left( X,d \right)$ is a metric space, we may consider the space $\ell^{\infty}\left( X \right)$ to be the space of all bounded sequences in $X$ with metric defined by
\begin{align*}
  d_{\infty}\left( x,y \right) &= \sup_{n\in \N} d\left( x_n,y_n \right).
\end{align*}
If $\omega$ is an ultrafilter on $\N$, then we define an equivalence relation $x\sim y$ on $\ell^{\infty}\left( X \right)$ by saying they are equal if
\begin{align*}
  \lim_{\omega} d\left( x_n,y_n \right) &= 0.
\end{align*}
Then, we define the equivalence classes
\begin{align*}
  \left[ x \right]_{\omega} &= \set{y\in \ell^{\infty}\left( X \right) | y\sim x}
\end{align*}
and take the quotient $X_{\omega} = \ell^{\infty}\left( X \right)/\sim$, with distance metric defined by
\begin{align*}
  d_{\omega}\left( \left[ x \right]_{\omega},\left[ y \right]_{\omega} \right) &= \lim_{\omega}d\left( x_n,y_n \right).
\end{align*}
This space is known as the \textit{ultrapower} of the metric space $\left( X,d \right)$. The map $\iota_{\omega}\colon X\rightarrow X_{\omega}$ defined by $x\mapsto \left[ \left( x,x,\dots \right) \right]_{\omega}$ is an isometric embedding of $\left( X,d \right)$ into $\left( X_{\omega},d_{\omega} \right)$.

Generally, we assume that the ultrafilter $\omega$ is nonprincipal, since if $\omega$ is principal, then taking the limit along the ultrafilter is equivalent to extracting the element of $\N$ that defines the principal ultrafilter.

A useful fact is that if $\left( X,d \right)$ is complete, then so too is $\left( X_{\omega},d_{\omega} \right)$.

Now, suppose $G$ is a group, and let $d\colon G\times G\rightarrow \left[ 0,\infty \right)$ be a metric on $G$. We say the metric $d$ is left (or right) invariant if $d\left( hg_1,hg_2 \right) = d\left( g_1,g_2 \right)$ for left invariance and $d\left( g_1h,g_2h \right) = d\left( g_1,g_2 \right)$ for right invariance. A metric that is both left and right invariant is called bi-invariant, and the space $\left( G,d \right)$ is called a metric group.
\begin{example}
  Let $G$ be a finitely generated group, and let $S\subseteq G$ be a finite generating subset. Then, $\left( G,d_S \right)$ is a left metric group, where $d_S$ denotes the word metric --- i.e., $d_S\left( g,h \right)$ is the minimum word length necessary to construct the element $gh^{-1}$ from $S$.
\end{example}
In the next section, we will discuss two important metric groups and various important properties thereof.
\section{Sofic Groups}%
Suppose $S_n$ is the group of permutations of the set $[n] = \set{0,\dots,n-1}$. Then, the \textit{Hamming length} of a permutation $\sigma\in S_n$ is given by
\begin{align*}
  \ell \left( \sigma \right) &= \frac{1}{n}\left\vert \set{i\in [n] : \sigma\left( i \right)\neq i} \right\vert.
\end{align*}
The metric $d\colon S_n\times S_n\rightarrow [0,\infty)$ given by $d\left( x,y \right) = \ell\left( xy^{-1} \right)$ then is a bi-invariant length function.\footnote{For invariance, recall that conjugation preserves cycle type, which is essentially what the Hamming distance is measuring. One can show that a metric being bi-invariant is equivalent to it being preserved under conjugation.}
\begin{definition}
  A countable discrete group $\Gamma$ is \textit{sofic} if for every $\ve > 0$ and every finite subset $F$ of $\Gamma\setminus \set{e}$, there is $n$ and a map $\Phi\colon \Gamma\rightarrow S_n$ such that $\Phi(e) = \id$, and for every $g,h\in F\setminus \set{e}$, we have
  \begin{itemize}
    \item $d\left( \Phi\left( gh \right),\Phi\left( g \right)\Phi\left( h \right) \right) < \ve$;
    \item $\ell\left( \Phi\left( g \right) \right) > r(g)$, where $r(g)$ is a constant depending only on $g$.
  \end{itemize}
\end{definition}
One of the main benefits of the ultrapower method is that it allows for a conversion from this local definition to a global definition. Fixing a free ultrafilter $\omega$ on $\N$, we define the ultraproduct
\begin{align*}
  \prod_{\omega} S_n &\coloneq \prod_{n\in \N}S_n/\set{x\in \prod_{n\in \N}S_n : \lim_{\omega}\ell_{S_n}\left( \sigma \right) = 0}.
\end{align*}
It follows from this definition that $\Gamma$ is sofic if and only if there exists an injective homomorphism $\Phi\colon \Gamma\rightarrow \prod_{\omega}S_n$ for any free ultrafilter $\omega$ on $\N$.
\begin{definition}
  Suppose $G$ and $H$ are groups with invariant length functions, $F$ a subset of $G$, and $\ve$ is a positive number. A function $\Phi\colon G\rightarrow H$ is called an $\left( F,\ve \right)$-approximate homomorphism if $\Phi\left( e_G \right) = e_H$ and for every $g,h\in F$,
  \begin{itemize}
    \item $\left\vert \ell_H\left( \Phi\left( gh \right)\Phi\left( h \right)^{-1}\Phi\left( g \right)^{-1} \right) \right\vert < \ve$;
     \item $\left\vert \ell_H\left( \Phi(g) \right) - \ell_G(g) \right\vert < \ve$.
  \end{itemize}
\end{definition}
\begin{theorem}
  A countable discrete group $\Gamma$ is sofic if and only if for every $\ve > 0$ and every finite $F\subseteq \Gamma$, there is $n$ and an $\left( F,\ve \right)$-approximate homomorphism from $\Gamma$ to $S_n$, where $\Gamma$ is regarded as an invariant length group with respect to the trivial invariant length function.
\end{theorem}
Here, the trivial invariant length function refers to the function
\begin{align*}
  \ell_0(x) &= \begin{cases}
    0 & x = e_{\Gamma}\\
    1 & \text{else}.
  \end{cases}
\end{align*}
\subsection{Local Embeddability}%
We discuss a more general notion than soficity that we will specialize to the sofic case.
\begin{definition}
  Let $G$ and $C$ be groups. If $K\subseteq G$ is finite, a map $\varphi\colon G\rightarrow C$ is called a $K$-almost homomorphism of $G$ into $C$ if it satisfies
  \begin{itemize}
    \item $\varphi\left( k_1k_2 \right)= \varphi\left( k_1 \right)\varphi\left( k_2 \right)$ for all $k_1,k_2\in K$;
    \item $\varphi|_{K}\colon K\rightarrow C$ is injective.
  \end{itemize}
\end{definition}
We call a collection $\mathcal{C}$ a \textit{class of groups} if it is a collection of groups such that if $C\in \mathcal{C}$ is a member, and $C'\cong C$ are isomorphic, then $C'\in \mathcal{C}$. The main ones we will be interested in are the class of finite groups and the class of amenable groups.
\begin{definition}
  If $\mathcal{C}$ is a class of groups, then we say $G$ is \textit{locally embeddable} into $\mathcal{C}$ if for every finite subset $K\subseteq G$, there exists a group $C\in \mathcal{C}$ and a $K$-almost homomorphism $\varphi\colon G\rightarrow C$.
\end{definition}
\begin{proposition}
  Let $\mathcal{C}$ be a class of groups. Every subgroup of a group that is locally embeddable into $\mathcal{C}$ is locally embeddable into $\mathcal{C}$.
\end{proposition}
\begin{proof}
  Let $G$ be a group that is locally embeddable into $\mathcal{C}$, and let $H$ be a subgroup of $G$. Given a finite subset $K\subseteq H$, then $\varphi\colon G\rightarrow C$ is a $K$-almost homomorphism into $C\in \mathcal{C}$ such that $\varphi|_{H}\colon H\rightarrow C$ is a $K$-almost homomorphism of $H$ into $C$, so $H$ is locally embeddable into $\mathcal{C}$.
\end{proof}
\begin{proposition}
  Let $G$ be a group, and let $\mathcal{C}$ be a class of groups. Then, $G$ is locally embeddable into $\mathcal{C}$ if and only if every finitely generated subgroup of $G$ is locally embeddable into $\mathcal{C}$.
\end{proposition}
\begin{proof}
  One direction follows from the previous proposition.

  Now, suppose every finitely generated subgroup of $G$ is locally embeddable into $\mathcal{C}$. Let $K\subseteq G$ be a finite subset, and let $H = \left\langle K \right\rangle$ be the subgroup generated by $K$. Then, since $H$ is locally embeddable into $\mathcal{C}$, there is a $K$-almost homomorphism $\varphi\colon H\rightarrow C$ of $H$ into $C\in \mathcal{C}$. We may extend $\varphi$ to $G$ arbitrarily by setting $\varphi(g) = e_C$ for all $g\in G\setminus H$, whence $\varphi\colon G\rightarrow C$ is a $K$-almost homomorphism of $G$ into $C$.
\end{proof}
\subsection{Examples and Inheritance Properties}%
\begin{proposition}
  Every finite group is sofic.
\end{proposition}
\begin{proof}
  Let $G$ be a finite group. By Cayley's Theorem, we have a map $L\colon G\rightarrow \sym(G)$ defined by $L(g)(h) = gh$ for all $g,h\in G$. Since $L$ is a homomorphism, we have that $\ell_G\left( L\left( g_1g_2 \right)L\left( g_2 \right)^{-1}L\left( g_1 \right)^{-1} \right) = 0$. Similarly, since $\ell_G\left(L(g)\right) = \ell_G\left( g \right)$, it follows that this is a $\left( G,\ve \right)$-approximate homomorphism.
\end{proof}
\begin{proposition}
  Amenable groups are sofic.
\end{proposition}
\begin{proof}
  Let $\Gamma$ be a amenable countable discrete group, $F$ a finite subset of $\Gamma$, and $\ve$ a positive real number. Find a $\left( F\cup F^{-1},\ve \right)$-invariant subset $K$ of $\Gamma$, whose existence follows from the Følner condition. If $\gamma\in F$, then define
  \begin{align*}
    \sigma_{\gamma}(x) &= \gamma x
  \end{align*}
  for each $x\in \gamma^{-1}K\cap K$, and extend $\sigma_{\gamma}$ to an arbitrary permutation of $K$. Then, this defines a $\left( F,2\ve \right)$-approximate homomorphism to the group of permutations of $K$ endowed with the Hamming length. Therefore, $\Gamma$ is sofic.
\end{proof}
\section{Hyperlinear Groups}%

\nocite{topics_in_groups_and_geometry,cellular_automata_and_groups,capraro2015introductionsofichyperlineargroups}
\printbibliography
\end{document}
